% Build in this directory: pdflatex -interaction=nonstopmode -halt-on-error class7_handout.tex
\documentclass[a4paper,10pt,pdflatex,ja=standard,jafont=ipaex-type1,
  magstyle=real,scale=1]{bxjsarticle}
\usepackage{lmodern}
\usepackage{amsmath,amssymb,mathtools}
\usepackage{geometry}
\geometry{reset,left=22mm,right=22mm,top=20mm,bottom=17mm,
  headheight=13pt,headsep=6mm,footskip=10mm}
\usepackage{xcolor}
\usepackage{enumitem}
\usepackage{array,booktabs,tabularx}
\usepackage{fancyhdr}

\definecolor{ink}{HTML}{193A5C}
\definecolor{accent}{HTML}{237D82}
\definecolor{soft}{HTML}{EAF3F3}
\definecolor{muted}{HTML}{52616B}
\newcommand{\coursename}{解析学 II}
\newcommand{\handoutnumber}{7}
\pagestyle{fancy}
\fancyhf{}
\fancyhead[L]{\small\color{ink}\coursename{}・授業要約と演習}
\fancyhead[R]{\small\color{ink}第\handoutnumber 回}
\fancyfoot[C]{\small\color{ink}\thepage/2}
\renewcommand{\headrulewidth}{0.3pt}
\setlength{\parindent}{0pt}
\setlength{\parskip}{2pt}
\setlength{\abovedisplayskip}{5pt}
\setlength{\belowdisplayskip}{5pt}
\setlength{\abovedisplayshortskip}{3pt}
\setlength{\belowdisplayshortskip}{4pt}
\setlist[itemize]{leftmargin=1.8em,itemsep=1pt,topsep=2pt,parsep=0pt}

\newcommand{\handouttitle}[1]{%
  \begin{center}
    {\Large\color{ink}#1}\par\vspace{3mm}
    {\color{accent}\rule{0.88\linewidth}{0.7pt}}
  \end{center}\vspace{-2mm}}
\newcommand{\handout}[2]{%
  \renewcommand{\handoutnumber}{#1}%
  \handouttitle{第#1回\quad #2}}
\newcommand{\topic}[1]{\par\vspace{5pt}%
  {\large\sffamily\mdseries\color{accent}#1}\par\vspace{2pt}}
\newcommand{\keybox}[1]{\par\vspace{3pt}\begingroup
  \setlength{\fboxsep}{4pt}%
  \colorbox{soft}{\parbox{\dimexpr\linewidth-2\fboxsep\relax}{\small #1}}%
  \endgroup\par\vspace{2pt}}
\newenvironment{problems}
  {\begin{enumerate}[leftmargin=2em,itemsep=3pt,topsep=3pt,parsep=0pt,
    font=\color{accent}]}
  {\end{enumerate}}


\newcommand{\examplelabel}[1]{{\sffamily\mdseries\color{ink}#1}}
\newcommand{\answerroom}{\par\vspace{16mm}}

\begin{document}
\fontsize{10pt}{13.5pt}\selectfont
\setlength{\abovedisplayskip}{4pt}
\setlength{\belowdisplayskip}{4pt}
\handout{7}{重積分の応用：体積と質量}
\topic{上下のグラフに挟まれた立体の体積}
$D\subset\mathbb R^2$ を有界な積分可能領域、$\phi,\psi$ を連続関数、$\phi\le\psi$ とし
$$
 V=\{(x,y,z):(x,y)\in D,\ \phi(x,y)\le z\le\psi(x,y)\}
$$
とする。各点の柱の高さは $\psi-\phi$ なので
$$
 \operatorname{Vol}(V)=\iint_D(\psi-\phi)\,dx\,dy
 =\iiint_V1\,dx\,dy\,dz.
$$
上面と下面を入れ替えず、投影領域 $D$ と高さの範囲を先に決める。

\topic{例：半径 $R>0$ の球の体積}
$D=\{(x,y):x^2+y^2\le R^2\}$ の上で
$\psi=\sqrt{R^2-x^2-y^2}$、$\phi=-\psi$ とすると
$$
 \begin{aligned}
 \operatorname{Vol}(V)
 &=\int_0^{2\pi}\int_0^R2\sqrt{R^2-r^2}\,r\,dr\,d\theta\\
 &=4\pi\left[-\frac13(R^2-r^2)^{3/2}\right]_0^R
 =\frac{4\pi R^3}{3}.
 \end{aligned}
$$

\topic{密度を積分して総量を求める}
質量密度 $\rho\ge0$ が与えられ、$\rho$ が $V$ 上可積分なら、質量は
$$
 M=\iiint_V\rho(x,y,z)\,dx\,dy\,dz
 =\iint_D\left(\int_{\phi(x,y)}^{\psi(x,y)}\rho(x,y,z)\,dz\right)dx\,dy.
$$
密度が一定 $\rho_0$ なら $M=\rho_0\operatorname{Vol}(V)$。
電荷密度・エネルギー密度を積分すれば、それぞれ全電荷・全エネルギーになる。

\topic{円柱座標と球面座標の体積要素}
円柱座標 $x=r\cos\theta,\ y=r\sin\theta$ では
$$
 dx\,dy\,dz=r\,dr\,d\theta\,dz.
$$
球面座標を $x=s\sin\alpha\cos\theta,\ y=s\sin\alpha\sin\theta,\ z=s\cos\alpha$
（$s\ge0,\ 0\le\alpha\le\pi$）と定めると
$$
 dx\,dy\,dz=s^2\sin\alpha\,ds\,d\alpha\,d\theta.
$$
円柱・球の対称性に合う座標を選ぶと、領域と被積分関数が簡単になる。

\topic{例：円柱の質量と重心}
$V=\{x^2+y^2\le1,\ 0\le z\le2\}$、$\rho=z$ とする。
$$
 M=\int_0^{2\pi}\int_0^1\int_0^2z r\,dz\,dr\,d\theta=2\pi.
$$
一般に $M>0$ の物体の重心は $\overline x=M^{-1}\iiint_Vx\rho\,dV$ などで定まる。
この例では回転対称性から $\overline x=\overline y=0$、また
$$
 \overline z=\frac1{2\pi}\int_0^{2\pi}\int_0^1\int_0^2z^2r\,dz\,dr\,d\theta=\frac43.
$$
\keybox{体積は1を、質量は密度を積分する。座標を変換したら、体積要素のヤコビアンも必ず付ける。}

\newpage
\handouttitle{第7回\quad 演習問題}
計算過程を記し、必要な条件・積分領域・向きを明示すること。
\begin{problems}
\item $V=\{(x,y,z):x^2+y^2\le1,\ 0\le z\le1-x^2-y^2\}$ の体積を求めよ。投影領域と上下の境界を明記すること。\answerroom
\item 半径 $R>0$ の球の体積を球面座標で求め、上下のグラフを用いた計算と一致することを確認せよ。\answerroom
\item 半径1・高さ2の円柱に一定密度 $\rho_0>0$ があるときの質量を求めよ。次に密度を $\rho=z$ に変えた場合の質量を求めよ。\answerroom
\item $V=\{(x,y,z):x,y,z\ge0,\ x+y+z\le1\}$ の体積を累次積分で求めよ。\answerroom
\item $V=\{(x,y,z):0\le x,y,z\le1\}$、$\rho=1+z$ とする。質量 $M$ と重心の高さ $\overline z$ を求めよ。\answerroom
\item 半径 $R>0$ の球内で $\rho(x,y,z)=x^2+y^2+z^2$ とする。質量を球面座標で求め、重心が原点となる理由を述べよ。\answerroom
\end{problems}
\end{document}
